\section{Network definition and EOC parameterization}\label{sec:network-definition}

\vspace{1em}
\subsection{RF-LR architecture}
Let \( \bm{h}^{(0)}(x)=x\in\mathbb{R}^{\bm{d}_0} \). For \( \ell=1,\dots,L \),
\begin{equation}
\label{eq:RF-LR}
\bm{h}^{(\ell)}(x)
= \frac{1}{\sqrt{\bm{n}_\ell}} \sum_{j=1}^{\bm{n}_\ell} A^{(\ell)}_j\,
\bm{\ReLU}\!\Big(w^{(\ell)\top}_j\,\bm{h}^{(\ell-1)}(x) + b^{(\ell)}_j\Big) \;+\; c.
\end{equation}
\textbf{Training policy:} \( A^{(\ell)} \) and the scalar output bias \( c \) are trained; \( w^{(\ell)} \) and activation biases \( b^{(\ell)} \) are frozen random draws (i.i.d. Gaussian). The \( 1/\sqrt{\bm{n}_\ell} \) scaling (NTK parameterization) ensures a well-defined infinite-width limit.
\paragraph{Width vs bottleneck dimension.}
We distinguish the feature width \(\bm{n}_\ell\) (number of random features) from the output dimension \(\bm{d}_\ell=\dim(\bm{h}^{(\ell)})\). The trainable matrix \(A^{(\ell)}\in\mathbb{R}^{\bm{d}_\ell\times \bm{n}_\ell}\) has columns \(A^{(\ell)}_j\in\mathbb{R}^{\bm{d}_\ell}\). A \emph{bottleneck} layer corresponds to \(\bm{d}_\ell=r\ll \bm{n}_\ell\), yielding \(O(r\bm{n}_\ell)\) trainable parameters and \(\mathrm{rank}(A^{(\ell)})\le r\) automatically.

\paragraph{Training parameters, biases, and initialization.}
Initialization is
\[
w^{(\ell)}_j \sim \mathcal{N}\!\big(0, I_{\bm{d}_{\ell-1}}/\bm{d}_{\ell-1}\big),\quad
b^{(\ell)}_j \sim \mathcal{N}(0, \beta^2),\quad
A^{(\ell)}_{ij} \sim \mathcal{N}\!(0, \sigma_A^2),\quad
c \sim \mathcal{N}(0, \sigma_c^2),
\]
We use the \emph{edge-of-chaos (EOC)} scaling for the weights so that \( \mathrm{Cov}\!\big(w^{(\ell)}_j\big) = I_{\bm{d}_{\ell-1}}/\bm{d}_{\ell-1} \) \cite{hayou2018selectioninitializationactivationfunction}.
In our bottleneck regime we take \(\bm{d}_{\ell-1}=r\) for \(\ell\ge 2\), so \(w^{(\ell)}_j\in\mathbb{R}^r\) with covariance \(I_r/r\).

We include an activation bias \( b^{(\ell)} \) and a scalar output bias \( c \) for completeness, but we avoid introducing extra scaling parameters for them. Output scaling can be absorbed into the learning rate; the mean-field and NTK structures remain qualitatively identical, and the Fourier-space interpretation and stepwise loss phenomenon depend on kernel geometry, not absolute scales.

With this setup in place, we now recall a minimal assumption on the activation that we will use throughout.

\paragraph{Activation homogeneity.}\label{assump:homogeneous_activation}
We only assume a positively 1-homogeneous nonlinearity \(\bm{\ReLU}\) for Fisher–Kibble decoupling (see Lemma~\ref{lem:fisher-kibble}):
\[
\ReLU(\alpha u) \,=\, \alpha\,\ReLU(u) \quad \text{for } \alpha \ge 0,\; u \in \mathbb{R}.
\]
This yields the scaling relation \(\ReLU(\alpha\, w^\top h) = \alpha\,\ReLU(w^\top h)\) and, away from zero, \(\dot{\ReLU}(\alpha u)=\dot{\ReLU}(u)\).

\subsection{Parameterization, signal propagation, and kernels}

With the model and initialization fixed, we can now examine parameterization choices, signal propagation at initialization, and the kernel quantities driving our analysis. We begin by tracking the layer-wise variance at initialization.

For general weight variance \(\Sigma_w = I_{\bm{d}_{\ell-1}}/\bm{d}_{\ell-1}\), define the per-layer variance \( q^{(\ell)} \). By weight independence we obtain the recursion:
\[
q^{(\ell)} = \sigma_c^2 + \frac{\sigma_A^2}{2} q^{(\ell-1)}\,\mathrm{Tr}(\Sigma_w), \]
This leads to the \emph{(EOC)} condition: 
\[ \frac{\sigma_A^2}{2}\mathrm{Tr}(\Sigma_w)=1. \]

It is chosen to prevent forward activations and backward sensitivities from exploding or vanishing across depth. For example, under \(\Sigma_w = I_{\bm{d}_{\ell-1}}/\bm{d}_{\ell-1}\), one obtains \(\sigma_A = \sqrt{2}\) (and we keep \(\sigma_c^2\) fixed).

\paragraph{Origin of \( 1/r \) in the NTK recursion.}
For bottleneck layers \( \ell \ge 2 \), the input dimension is \( \bm{d}_{\ell-1} = r \), so EOC prescribes \( \Sigma_w = I_r/r \) with \( \mathrm{Tr}(\Sigma_w) = 1 \). The derivative kernel
\[
\dot{\Sigma}^{(\ell)}(x,x') = \mathbb{E}_w\big[\dot{\sigma}(w^\top h)\,\dot{\sigma}(w^\top h')\,\|w\|^2\big] = \mathrm{Tr}(\Sigma_w)\!\left(\tfrac{1}{2} - \tfrac{\arccos(\rho_\ell)}{2\pi}\right)
\]
remains \( O(1) \) under EOC. Moreover, with \(\mathrm{Tr}(\Sigma_w)=1\) this quantity is uniformly bounded as
\[
0\le \dot{\Sigma}^{(\ell)}(\rho)\le \frac12,\qquad \forall \rho\in[-1,1],
\]
which can be used as an explicit choice of the constant \(c_0\) in large-depth bounds based on products of derivative kernels.
In our RF-LR setting the NTK involves gradients with respect to the bottleneck readouts \(A^{(\ell)}\), and the resulting recursion carries an explicit \(1/r\) prefactor in the bottleneck regime (Theorem~\ref{thm:ntk_recursion}). Intuitively, this factor reflects the bottleneck scaling in the parameterization and the fact that the relevant Jacobian contributions concentrate in an \(r\)-dimensional subspace. We provide the full derivation in Appendix~\ref{appendix:proof-ntk-recursion}.
\[
\Theta^{(\ell)} = 1 + \frac{1}{r}\,\Theta^{(\ell-1)}\cdot\sigma_A^2\,\dot{\Sigma}^{(\ell)} + \frac{1}{r}\,\sigma_A^2\,\Sigma^{(\ell)}, \qquad \ell \ge 2.
\]
In summary: EOC fixes \( \mathrm{Tr}(\Sigma_w)=1 \) so that \( \dot{\Sigma}^{(\ell)}, \Sigma^{(\ell)} \) are \( O(1) \), while the low-rank/bottleneck scaling appears through the explicit prefactor \(1/r\) in the NTK recursion.

\vspace{0.5em}
With an abuse of notation, we denote by \( \Sigma^{(\ell)} \) the base kernel and by \( \dot{\Sigma}^{(\ell)} \) the derivative kernel. We emphasize that these are not the NNGP kernel and the derivative kernel, but the base kernel and the derivative kernel of the Gaussian process \( \bm{h}^{(\ell)} \) conditional on \( \bm{h}^{(\ell-1)} \).
By rotation invariance of Gaussian weights, these kernels depend on inputs only through the cosine similarity \(\rho=\langle x,x'\rangle/(\|x\|\,\|x'\|)\). We will therefore write \(\Sigma^{(\ell)}(\rho)\) and \(\dot{\Sigma}^{(\ell)}(\rho)\). When no confusion arises (e.g., for \(\ell=1\)), we will also abuse notation and write \(\sigma(\rho)\) for the scalar kernel function; this \(\sigma\) denotes the kernel-as-a-function-of-\(\rho\), not the activation \(\bm{\ReLU}\).
We can now state the base and derivative kernels precisely.
\begin{definition}[Base and derivative kernels]\label{def:base-derivative-kernels}
\vspace{0.3cm}
In the sequential infinite-width limit, the layer-\( \ell \) base kernel is
\begin{equation}
\Sigma^{(\ell)}(x,x')
= \mathbb{E}_{w}\Big[\bm{\ReLU}\!\big(w^\top \bm{h}^{(\ell-1)}(x)\big)\,\bm{\ReLU}\!\big(w^\top \bm{h}^{(\ell-1)}(x')\big)\Big],
\end{equation}
and the derivative kernel is
\begin{equation}
\dot{\Sigma}^{(\ell)}(x,x')
= \mathbb{E}_{w}\Big[\dot{\bm{\ReLU}}\!\big(w^\top \bm{h}^{(\ell-1)}(x)\big)\,\dot{\bm{\ReLU}}\!\big(w^\top \bm{h}^{(\ell-1)}(x')\big)\;\|w\|^2\Big].
\end{equation}
For \( \ell=1 \) and centered Gaussian \( w \), \( \Sigma^{(1)} \) is rotation-invariant with standard closed forms in terms of input cosine similarity \( \rho = \langle x,x'\rangle/(\|x\|\|x'\|) \):
\[
\Sigma^{(1)}(x,x') = \frac{\|x\|\|x'\|}{2\pi}\big((\pi-\theta)\cos\theta + \sin\theta\big), \quad \theta = \arccos(\rho).
\]
\end{definition}

Given these definitions, we can now state the following composition result.
\begin{theorem}[Gaussian process composition]\label{thm:gp-composition}
\vspace{0.3cm}
    In the sequential infinite-width limit, the hidden states \( \bm{h}^{(\ell)} \) form a composition of Gaussian processes:
\begin{itemize}
    \item Conditionally on \( \bm{h}^{(\ell-1)} \), for any finite collection of inputs \( \{x_1, \dots, x_m\} \), the vector \( (\bm{h}^{(\ell)}(x_1), \dots, \bm{h}^{(\ell)}(x_m)) \) converges in distribution to a multivariate Gaussian.
    \item The conditional limiting Gaussian has mean zero and covariance matrix \( [\Sigma^{(\ell)}(x_i, x_j)]_{i,j=1}^m \), where \( \Sigma^{(\ell)} \) is the base kernel defined in Definition~\ref{def:base-derivative-kernels} and depends on \( \bm{h}^{(\ell-1)} \)
\end{itemize}
\end{theorem}
\begin{proof}
See Appendix~\ref{appendix:proof-gp-composition}.
\end{proof}


\paragraph{Composition of Gaussian processes} We highlight the main difference between the base kernel and the NNGP kernel arising in classical NTK analysis : Since \( \bm{h}^{(\ell-1)} \) itself is a (conditional) Gaussian process, the overall distribution is a composition of Gaussian processes, not a single Gaussian process.

